% DO -- os/samir/cmd/proc.c proc_load_program / p_do / p_call (one loader
% for the dot prompt and for a program; a SET PROCEDURE file wins over a
% same-named disk program. Current-source precedence is not captured. 8191-byte
% cap, #18; WITH arguments by value).  Real III+ 1.1 captures:
% xbase_prog_diff/do-file-mint/{basic,depth,handler,master,missing,
% parameters,precedence,close} (initech-104j3).  Gates: test-interp-do-file,
% test-samir-canon-y2k, test-interp-proc.
\begin{XbaseEntry}{DO}{command}
Runs a program file, or a procedure of the program that is running.
\Syntax
\SyntaxLine{\kw{DO} \xp{program}}
\par or\par
\SyntaxLine{\kw{DO} \xp{procedure} \opt{\kw{WITH} \xp{expression list}}}
\Usage
At the dot prompt, \cmd{DO} \xp{program} reads the program file
\xp{program}\texttt{.PRG} from the disk and runs it.  You need not type the
extension.  The program runs from its first line until it reaches
\cmd{RETURN}, its first \cmd{PROCEDURE} or its last line; then the dot prompt
returns.  A program file may hold at most 8,191 characters; a larger file is
refused with \xmsg{18}{Line exceeds maximum of 254 characters.}

In a program, \cmd{DO} \xp{name} can run a procedure in the open procedure
file (see \cmd{SET PROCEDURE}), a procedure in the same program file, or
the program file \xp{name}\texttt{.PRG} on the disk.  A procedure in the
open procedure file takes precedence over a disk program of the same name.
When it returns, the program goes on with the line
after the \cmd{DO}.  Procedures and programs may call others, and themselves.  At
most 20 \cmd{DO} commands may be in progress at once, counting the program
started at the dot prompt; the \cmd{DO} that would begin a twenty-first is
refused with \xmsg{103}{DOs nested too deep.}

\kw{WITH} passes the values of the expressions to the procedure, which
receives them in the memory variables of its \cmd{PARAMETERS} line, in order.
The procedure receives a copy of each value: a change it makes to a parameter
does not change the expression that was passed.
% Real III+ do-file-mint/parameters/{MINT.PRG,GOLD.OUT}: 'alpha' is passed
% literally, received intact, and remains intact after a nested disk DO.
A character constant may be written directly in the \kw{WITH} list.

In a program, a name found in none of those places is refused with
\xmsg{1}{File does not exist.}; at the dot prompt the refusal is
\xmsg{16}{*** Unrecognized command verb.}
\Tips
The programs that come with InitechBase work on the database already in USE:
open it at the dot prompt with \cmd{USE} before you type \cmd{DO}.
\Examples
\file{RAISE.PRG} calls its procedure \texttt{NEWSAL} with two values:
\ProgramListing{RAISE.PRG}
\DotExample{examples/c-do}
\SeeAlso{PARAMETERS, PROCEDURE, RETURN, PUBLIC, PRIVATE}
\end{XbaseEntry}

